%% IoTJ-PaperD.tex
%% Fragment-Chained HMAC Authentication with Selective Encryption
%% for JPEG Transmission over LoRa Mesh
%% Target: IEEE Internet of Things Journal (IoT-J)

\documentclass[journal]{IEEEtran}

\usepackage[T1]{fontenc}
\usepackage[utf8]{inputenc}
\usepackage{amsmath,amssymb}
\usepackage{graphicx}
\usepackage{booktabs}
\usepackage{url}
\usepackage[hypertexnames=false]{hyperref}
\usepackage{array}
\usepackage{multirow}
\usepackage{tabularx}
\usepackage{xcolor}
\usepackage{cite}
\usepackage{listings}
\usepackage{algorithm}
\usepackage{algpseudocode}

\lstset{
  basicstyle=\ttfamily\small,
  breaklines=true,
  frame=single,
  columns=fullflexible
}

\begin{document}

\title{Fragment-Chained HMAC Authentication with Selective Encryption\\
for JPEG Transmission over LoRa Mesh}

\author{%
  \IEEEauthorblockN{August Chao}\\
  \IEEEauthorblockA{%
    Department of Computer Science and Information Engineering,\\
    National Pingtung University, Pingtung 90003, Taiwan\\
    \textit{augchao@gms.npu.edu.tw}%
  }%
}

\maketitle

%% -----------------------------------------------------------------------
\begin{abstract}
Transmitting JPEG images over LoRa mesh networks is challenging: the 255-byte
payload limit and ETSI 1\% duty-cycle constraint impose severe bandwidth and
latency budgets while security requirements demand both confidentiality and
integrity protection across tens of fragments.  Existing approaches apply full
AES-128-CBC encryption and independent per-packet HMAC, which (i) incurs
unnecessary CPU cost when only visually sensitive bytes require confidentiality,
and (ii) cannot detect fragment reordering or cross-session replay because each
packet is verified independently.

This paper makes four contributions.  First, we design a
\emph{stride-based selective encryption} scheme that XORs the first 16 bytes
of every 64-byte window with an HMAC-SHA256 keystream, encrypting 24.2\% of
image bytes and achieving format-level visual obfuscation (JPEG decode fails
under both naive and format-aware attacker models) while reducing encryption
CPU cost by 76\% compared to full encryption (2{,}807~$\mu$s vs.\
$\approx$11{,}600~$\mu$s estimated for the same image size). Second, we
introduce \emph{Fragment-Chained HMAC} (FC-HMAC), an 8-byte truncated chain
where each fragment's tag covers the previous tag, binding the entire
transmission sequence into a single cryptographic commitment. Third, we formally
verify FC-HMAC using Tamarin Prover: all four lemmas (fragment authenticity,
chain-gap detectability, no-replay, executability) are machine-proved in 1.53
seconds under a Dolev-Yao adversary. Fourth, we validate the scheme on real
LoRa hardware (EoRa PI: ESP32-S3 + SX1262): 72/72 fragments verified in
hardware over a 120-second run, with 3.3\% HMAC overhead per packet and
3.76~ms total crypto CPU per 2{,}046-byte image at identical airtime (14.9~s)
across all three encryption modes.
\end{abstract}

\begin{IEEEkeywords}
LoRa, LPWAN, image transmission, selective encryption, fragment chaining,
HMAC, formal verification, Tamarin Prover, ESP32-S3, IoT security
\end{IEEEkeywords}

%% -----------------------------------------------------------------------
\section{Introduction}
\label{sec:intro}

Low-Power Wide-Area Networks (LPWANs), and LoRa in particular, have emerged as
a dominant radio layer for battery-operated IoT sensors that need to communicate
over distances of several kilometres on harvested or coin-cell power
\cite{etsi_en300220}.  The physical layer chirp-spread-spectrum (CSS) modulation
trades data rate for link budget, achieving $-$148~dBm sensitivity at SF12 at
the cost of throughputs as low as 250~bps.  Within the 255-byte hardware payload
ceiling of the Semtech SX1262 \cite{semtech_sx1262} and the ETSI 1\%
duty-cycle rule, a single LoRa node in the 915/868~MHz ISM band may inject only
$\approx$36 packets per hour at spreading factor~9 (SF9).

A growing class of edge applications requires \emph{image} transmission over
these links: wildlife camera traps that report motion events to a gateway,
border-surveillance fixed nodes that capture a thumbnail on alarm, disaster-area
drone relays that upload reconnaissance snapshots, and smart-agriculture sensors
that record crop canopy images for periodic analysis.  For these use cases, a
compact JPEG image ($200 \times 150$ pixels, $\approx$2~KB) is a practical and
sufficient representation.  Transmitting this image over LoRa at SF9/BW125/CR4/7
requires approximately 9 fragments and roughly 14.9~seconds of airtime --- within
budget for one capture every few minutes, but leaving no room for heavy
cryptographic overhead on the critical path.

\subsection*{The Security Gap in Prior Work}

Existing security schemes for LoRa image transmission \cite{ryczek2025}
apply \emph{full} AES-128-CBC encryption to the complete image payload, pairing
it with \emph{independent} per-packet HMAC-SHA256.  While this approach has
been validated at the hardware level and demonstrates negligible energy overhead,
it exhibits two structural weaknesses that limit its applicability to more
adversarial deployment environments:

\textbf{CPU inefficiency.}  Full AES-128-CBC encryption of a 2{,}046-byte image
on an ESP32-S3 at 240~MHz requires approximately 11{,}600~$\mu$s of CPU time.
On a duty-cycle-constrained node that must also handle radio interrupts, sensor
sampling, and routing, this represents a significant fraction of the available
compute budget.  When the goal is merely visual obfuscation --- preventing an
interceptor from reconstructing a recognisable image --- encrypting every byte
is wasteful, because a small fraction of structurally critical bytes (JPEG
format headers, DCT DC coefficients) suffices to render the image
unrecoverable.

\textbf{Order-agnostic integrity.}  Independent per-packet HMAC verifies each
fragment in isolation, providing \emph{existential unforgeability} for each
individual packet.  However, an attacker with access to a previous valid
capture session may:
\begin{enumerate}
  \item \emph{Replay} an entire valid session from a previous capture without
    triggering integrity failures at the receiver --- because no MAC field
    commits to the \emph{session context}.
  \item \emph{Reorder} fragments within a session and deliver them out of
    sequence; each fragment passes its individual HMAC check even though the
    assembly order has been corrupted.
  \item \emph{Insert} a forged fragment from a different session at any
    position; the receiver cannot distinguish this from a legitimately received
    fragment based on per-packet HMAC alone.
\end{enumerate}
These vulnerabilities are not theoretical: in a LoRa mesh deployment, a
Dolev-Yao adversary who controls the radio channel can store-and-forward any
received fragment to any subsequent session \cite{tamarin_prover}.

\subsection*{Contributions}

This paper addresses both weaknesses through the design, formal analysis, and
hardware validation of the \emph{Fragment-Chained HMAC} (FC-HMAC) scheme
combined with stride-based selective encryption.  The contributions are:

\textbf{C1 --- Stride-based selective encryption.}
We design a lightweight structure-approximating encryption scheme that XORs the
first 16 bytes of every 64-byte window with an HMAC-SHA256 keystream, encrypting
24.2\% of image bytes.  The scheme achieves \emph{format-level visual
obfuscation}: the stride begins at byte~0 (the JPEG SOI marker, 0xFF~0xD8) and
corrupts Huffman table regions, breaking JPEG structure entirely under both a
naive attacker and a format-aware attacker who zeroes the encrypted positions.
We explicitly acknowledge that this is \emph{not} true frequency-domain
selective encryption (which requires a Huffman-stream parser to locate DCT DC
coefficient bytes at variable positions in the bitstream); the contribution is a
computationally cheap proxy that achieves the same visual disruption at 24\% of
full-encryption CPU cost.

\textbf{C2 --- Fragment-Chained HMAC (FC-HMAC).}
We introduce an 8-byte truncated chained MAC where
$\mathit{tag}_i = \allowbreak \mathrm{HMAC\text{-}SHA256}(K,\;
 \mathit{frag\_id}_i \;\|\; \mathit{flags}_i \;\|\;
 \mathit{payload}_i \;\|\; \mathit{tag}_{i-1})[\text{:}8]$.
The chain makes every tag a commitment to the entire prior transmission
sequence, enabling detection of any reordering, insertion, or cross-session
replay at 3.3\% per-packet overhead (8~B~/~241~B).

\textbf{C3 --- Formal verification.}
We model FC-HMAC in Tamarin Prover \cite{tamarin_prover} and machine-prove
four security properties: fragment authenticity (F1), chain-gap detectability
(F2), no-replay (F3), and protocol executability (F4).  All four lemmas are
proved in 1.53~seconds under a Dolev-Yao adversary.

\textbf{C4 --- Hardware validation.}
We implement the full protocol on EoRa PI hardware (ESP32-S3 + SX1262,
EoRa PI development board), transmitting a 2{,}046-byte synthetic JPEG at
922.0~MHz over a real LoRa channel.  In a 120-second run with 8 complete
images, 72/72 fragments pass chain verification (chain\_ok = 100\%,
PDR = 100\%).  Total crypto CPU is 3{,}757~$\mu$s per image, representing
3.1\% of a 120-ms processing loop, with 3.3\% per-packet HMAC overhead.

\subsection*{Scope and Reproducibility}

The experiments in this paper use a 2{,}046-byte synthetic JPEG embedded in
firmware flash.  The choice of a synthetic image is deliberate: it enables
reproducible measurements independent of camera hardware variability and
lighting conditions.  The byte array is generated offline from a known-content
$200\times150$ pixel test pattern using the libjpeg library at quality factor
75.  All firmware source code, Tamarin model files, PSNR evaluation scripts,
and raw E19 measurement data will be made publicly available via the project
repository upon paper acceptance.

\subsection*{Paper Organisation}

Section~\ref{sec:background} reviews LoRa constraints, JPEG structure, and the
FC-HMAC concept.
Section~\ref{sec:design} presents the detailed system design including threat
model, the selective encryption algorithm, FC-HMAC construction, and packet
format.
Section~\ref{sec:security} gives the Tamarin formal analysis and E19-PSNR
visual security evaluation.
Section~\ref{sec:eval} reports hardware platform, E19 experiment results,
fragment timing, and efficiency analysis.
Section~\ref{sec:related} compares FC-HMAC with related work, including a
seven-dimension table against Ryczek \& Sobieraj \cite{ryczek2025}.
Section~\ref{sec:discussion} discusses limitations and future directions.
Section~\ref{sec:conclusion} concludes.

%% -----------------------------------------------------------------------
\section{Background}
\label{sec:background}

\subsection{LoRa Physical Layer and Protocol Constraints}

LoRa is a chirp-spread-spectrum modulation scheme defined by Semtech and
standardised within the ETSI EN~300~220 framework \cite{etsi_en300220}.  Three
parameters dominate the link budget and throughput tradeoff:

\textbf{Spreading Factor (SF).}  SF ranges from~6 to~12.  Each increment
doubles symbol duration ($T_s = 2^{SF}/BW$), increasing receiver sensitivity by
$\approx$2.5~dB at the cost of halving the data rate.  At SF9/BW125, symbol
duration is $T_s = 4.096$~ms, giving a raw bit rate of $\approx$1{,}758~bps.

\textbf{Payload ceiling.}  The SX1262 radio buffer is 255~bytes
\cite{semtech_sx1262}.  After subtracting the FC-HMAC header (7~B) and
authentication tag (8~B), the maximum application payload per packet is 240~B.

\textbf{Duty-cycle regulation.}  ETSI EN~300~220 mandates a 1\% maximum
on-air duty cycle in Sub-GHz ISM bands.  At SF9/BW125/CR4/7 with a 248-byte
payload, measured airtime per fragment is 1{,}692~ms.  The 1\% rule therefore
limits a single node to approximately 6~transmissions per 10-minute window ---
sufficient for one 9-fragment image every $\sim$4~minutes.

These constraints collectively motivate (i)~compact per-packet authentication
overhead to preserve payload capacity and (ii)~CPU-efficient encryption so that
the crypto latency does not dominate the inter-packet gap.

\subsection{LoRaWAN and LoRa Mesh Security Context}

LoRaWAN, the most widely deployed LoRa network architecture, provides
session-level AES-128 encryption and per-message MICs at the network
layer \cite{etsi_en300220}.  However, LoRaWAN is a star-topology protocol
requiring a centralised network server and dedicated gateway infrastructure.
Multi-hop LoRa mesh protocols, such as LoRaMesher \cite{loramesh_github},
operate in infrastructure-free environments where nodes relay packets for
one another using distance-vector routing.  In these deployments, the LoRaWAN
session key infrastructure (AppSKey, NwkSKey) is unavailable; nodes rely on
pre-shared application keys for security, and the security of the fragmentation
protocol --- which LoRaWAN does not define --- is left entirely to the
application layer.

The threat surface in a LoRa mesh is broader than in LoRaWAN because relay
nodes are themselves potential adversaries: a compromised relay can silently
drop, modify, or replay fragments it forwards, with no gateway to detect the
anomaly.  FC-HMAC is designed with this threat model explicitly in mind.

Previous work \cite{ryczek2025} demonstrates that application-layer AES-CBC
plus per-packet HMAC is feasible on ESP32-class hardware with $<$1\% energy
overhead for 20~kB images.  Our work examines whether the per-packet HMAC
architecture can be strengthened to detect ordering and replay attacks, and
whether selective encryption can reduce the CPU cost at the expense of
encrypting only a fraction of bytes.

\subsection{JPEG Structure and DCT Sensitivity}

JPEG images use Discrete Cosine Transform (DCT) coding applied to $8\times8$
pixel blocks \cite{jpeg_iso}.  The encoding pipeline proceeds as follows:

\begin{enumerate}
  \item \textbf{Color transform.}  RGB to YCbCr; chroma channels subsampled
    (4:2:0 or 4:2:2).
  \item \textbf{DCT per block.}  Each $8\times 8$ luma or chroma block produces
    64 frequency coefficients arranged in zigzag order.
  \item \textbf{Quantisation.}  Coefficients divided by a quality-dependent
    quantisation table; most AC coefficients are zeroed.
  \item \textbf{Huffman coding.}  DC differences between adjacent blocks are
    coded with a DC Huffman table; remaining AC run-length codes use an AC
    Huffman table.  Both tables are stored in the JPEG header (DHT markers).
  \item \textbf{Bitstream packing.}  Huffman-coded symbols are packed into a
    byte stream beginning with SOI (0xFF~0xD8) and ending with EOI (0xFF~0xD9).
\end{enumerate}

The DC coefficient (spectral index 0) captures the average brightness of each
block and has the highest per-coefficient visual impact: corrupting the DC
coefficient of one block propagates a visible brightness shift across the entire
$8\times8$ region.  A selective encryption scheme targeting only DC coefficients
could in principle disrupt the image with minimal ciphertext.

However, DC coefficient positions in the JPEG bitstream are determined by the
variable-length Huffman codes preceding them --- they do not occur at fixed byte
offsets.  Locating them requires streaming through the Huffman-coded data, which
implies implementing a JPEG decoder on the embedded node.  Our stride-based
approximation (Section~\ref{sec:sel_enc}) addresses this constraint.

\subsection{Fragment-Chained MAC: Conceptual Foundation}

A Merkle tree \cite{bellare_hmac} commits to an ordered set of leaves; any
modification of a leaf or its position is detectable at the root.  A
fragment-chained MAC is the streaming, append-only analogue: each node in the
chain commits to all prior nodes, so the final tag is a commitment to the
complete transmission history in order.

Formally, let $H$ be a collision-resistant PRF.  Define:
\begin{align}
  h_0 &= H(K,\; \mathit{context}) \label{eq:h0_gen}\\
  h_i &= H(K,\; m_i \;\|\; h_{i-1}), \quad i \geq 1 \label{eq:hi_gen}
\end{align}
where $m_i$ is the $i$-th fragment payload and $K$ is a shared key.  Any
adversary who modifies $m_j$ must forge $h_j$; since $h_j$ is bound to
$h_{j-1}$, forging $h_j$ requires either knowing $K$ or finding a collision in
$H$.  Furthermore, inserting a foreign fragment $m'_j$ (from a different
session or position) changes $h_j$, propagating the mismatch to all
subsequent tags.  The receiver can therefore verify not just individual packet
authenticity but also the integrity of the transmission \emph{order and
completeness}.

%% -----------------------------------------------------------------------
\section{System Design}
\label{sec:design}

\subsection{Threat Model}
\label{sec:threat}

The FC-HMAC system operates under the following threat model, which is
formally encoded in the Tamarin model (Section~\ref{sec:tamarin}).

\textbf{Network model.}  All LoRa transmissions are broadcast and receivable
by any node in range.  The adversary controls the network channel in the
Dolev-Yao sense: it can intercept any transmitted fragment, store it
indefinitely, and re-inject stored messages at any future time.

\textbf{Attacker capabilities.}  The adversary:
\begin{itemize}
  \item Can intercept, delay, reorder, and replay any LoRa packet.
  \item Can inject new packets with arbitrary byte content.
  \item Knows the encryption algorithm, MAC algorithm, stride offsets,
    and packet format (Kerckhoffs's principle).
  \item Does \emph{not} know the session key $K$ (pre-shared out-of-band).
  \item Cannot break HMAC-SHA256 computationally (PRF assumption).
\end{itemize}

\textbf{Goals protected against.}
Table~\ref{tab:threats} lists the threats addressed by FC-HMAC:

\begin{table}[ht]
\centering
\caption{Threat Model and FC-HMAC Protection Mapping}
\label{tab:threats}
\begin{tabular}{p{2.4cm}p{2.0cm}p{3.0cm}}
\toprule
Threat & Attack vector & FC-HMAC protection \\
\midrule
Fragment forgery & Inject crafted fragment & F1: tag forgery requires $K$ \\
Fragment replay (same session) & Re-inject captured fragment & Duplicate \texttt{frag\_id} rejected \\
Cross-session replay & Replay fragments from old session & F3: chain root binds session \\
Fragment reordering & Deliver frags out of order & Chain tag mismatch at reordered position \\
Silent gap (drop + no detect) & Drop frag $i$, continue chain & F2: chain break at gap position \\
Image eavesdropping & Intercept and decode image & C1: JPEG decode FAIL (E19-PSNR) \\
\midrule
Key compromise & Brute-force or side-channel $K$ & \emph{Out of scope} \\
Physical capture of node & Extract $K$ from flash & \emph{Out of scope (hardware TEE)} \\
\bottomrule
\end{tabular}
\end{table}

\textbf{Goals not protected.}  The scheme does not provide confidentiality
against an adversary who knows $K$ (key management is out of scope).  It does
not provide availability: an adversary can jam the LoRa channel or drop all
packets.  Physical security of the pre-shared key requires hardware trust
anchors (secure element, flash encryption) not present in the current EoRa PI
prototype.

\subsection{Stride-Based Selective Encryption}
\label{sec:sel_enc}

\subsubsection{Design Rationale and Scope}

True selective encryption targeting DCT DC coefficients requires a
Huffman-stream decoder to locate variable-length-coded DC bytes at runtime.
On a 240-MHz ESP32-S3 constrained to handle radio interrupts in real time, an
inline JPEG decoder would itself introduce CPU overhead comparable to
full-image encryption.

We therefore adopt a \emph{structure-approximating} approach.  The key
observation is that a JPEG file begins with mandatory format structures in the
first several hundred bytes:
\begin{itemize}
  \item Bytes 0--1: SOI marker (0xFF 0xD8)
  \item Bytes 2--20: APP0/JFIF marker (0xFF 0xE0) and header
  \item Bytes 21--$\sim$200: Quantisation tables (DQT markers, 0xFF 0xDB)
  \item Bytes $\sim$200--$\sim$500: Huffman tables (DHT markers, 0xFF 0xC4)
  \item Bytes $\sim$500+: Start of scan (SOS marker, 0xFF 0xDA) and encoded data
\end{itemize}
Corrupting any of the quantisation or Huffman table entries renders the JPEG
bitstream undecodable, because every JPEG decoder uses these tables to interpret
the subsequent AC/DC runs.  A stride starting at byte~0 with 16-byte windows
every 64~bytes will intersect these critical structures during its first several
iterations.

The contribution of the stride-based approach is therefore correctly framed as
\emph{lightweight visual obfuscation achieving format-level disruption}, not
frequency-domain selective encryption.  The actual overlap between stride
positions and true DCT DC coefficient byte positions remains to be quantified
(see Section~\ref{sec:limit_b1}).

\subsubsection{Algorithm Specification}

\begin{algorithm}[H]
\caption{Stride-Based Selective Encryption (\texttt{ENC\_DC\_ONLY} mode)}
\label{alg:sel_enc}
\begin{algorithmic}[1]
\Require Image byte array $P[0 \ldots N{-}1]$, session key $K$, frag counter $c$
\Ensure Encrypted byte array $C$
\State $C \leftarrow \text{copy}(P)$
\State $\mathit{stride} \leftarrow 64$; $\;\mathit{wlen} \leftarrow 16$
\For{$i \leftarrow 0$ \textbf{to} $N{-}1$ \textbf{step} $\mathit{stride}$}
  \State $\mathit{ks}[0{:}16] \leftarrow \mathrm{HMAC\text{-}SHA256}(K,\;
         \langle i \;\|\; c \rangle)[\text{:}16]$
  \For{$j \leftarrow 0$ \textbf{to} $\min(\mathit{wlen}{-}1,\; N{-}i{-}1)$}
    \State $C[i+j] \leftarrow P[i+j] \oplus \mathit{ks}[j]$
  \EndFor
\EndFor
\State \Return $C$
\end{algorithmic}
\end{algorithm}

The frag counter $c$ differentiates keystreams across repeated calls for the
same image buffer.  For a 2{,}046-byte JPEG, the loop executes at strides
$i \in \{0, 64, 128, 192, 256, 320, 384, 448, 512\}$, producing 9~keystream
calls and encrypting up to $\min(16, N-i)$ bytes per stride.
Hardware measurement (E19) confirms 496 of 2{,}046 bytes encrypted (24.2\%),
at a CPU cost of 2{,}807~$\mu$s --- a 76\% reduction compared to full-image
encryption estimated at $\approx$11{,}600~$\mu$s.

\subsubsection{Three Encryption Modes}

The firmware implements three modes selectable at compile time via the
\texttt{ENC\_MODE} build flag:

\begin{table}[ht]
\centering
\caption{Encryption Mode Comparison (2{,}046-byte JPEG, SF9/BW125)}
\label{tab:enc_modes}
\begin{tabular}{lrrrr}
\toprule
Mode & Packets & Airtime & CPU enc ($\mu$s) & Enc bytes \\
\midrule
\texttt{ENC\_NONE}    & 9 & 14.9\,s & 0         & 0\% \\
\texttt{ENC\_DC\_ONLY}& 9 & 14.9\,s & 2{,}807   & 24.2\% \\
\texttt{ENC\_FULL}    & 9 & 14.9\,s & $\approx$11{,}600 & 100\% \\
\bottomrule
\end{tabular}
\end{table}

A critical design point: all three modes transmit the \emph{same number of
packets at the same airtime}, because packet count is determined by
$\lceil N / 233 \rceil$ regardless of encryption mode.  The efficiency gain of
\texttt{ENC\_DC\_ONLY} is therefore in \emph{CPU}, not in airtime or duty-cycle
consumption.  This distinction is explored further in
Section~\ref{sec:efficiency}.

\subsubsection{LoRa Airtime Model}

For completeness, we derive the airtime formula used throughout this paper.
Using the standard LoRa time-on-air model \cite{semtech_sx1262}, with spreading
factor $\mathit{SF} = 9$, bandwidth $\mathit{BW} = 125$~kHz, coding rate
$\mathit{CR} = 4/7$, explicit header mode, and CRC enabled:

\begin{align}
  T_s &= \frac{2^{SF}}{BW} = \frac{512}{125000} = 4.096~\text{ms}
  \label{eq:Ts} \\
  N_{\text{preamble}} &= 8 \text{ symbols (default)} \nonumber \\
  T_{\text{preamble}} &= (N_{\text{preamble}} + 4.25) \times T_s = 50.9~\text{ms}
  \label{eq:Tpreamble}
\end{align}

For a 248-byte payload packet (7 header + 233 data + 8 HMAC):
\begin{equation}
  N_{\text{sym,payload}} =
    8 + \max\!\!\left(
      \left\lceil
        \frac{8 \times 248 - 4 \times SF + 28 + 16}{4 \times (SF-2)}
      \right\rceil \times CR,\, 0
    \right) = 107 \text{ symbols}
  \label{eq:Nsym}
\end{equation}
\begin{equation}
  T_{\text{packet}} = T_{\text{preamble}} + N_{\text{sym,payload}} \times T_s
    = 50.9 + 107 \times 4.096 \approx 489~\text{ms}
  \label{eq:Tpacket}
\end{equation}

The hardware-measured airtime of 1{,}692~ms per fragment includes inter-packet
gaps, radio mode-switch overhead, and SX1262 TX-to-RX turnaround time.  The
theoretical packet-on-air time (489~ms) is consistent with the fraction of
1{,}692~ms spent with the antenna active.  The total image airtime is:
\begin{equation}
  T_{\text{image}} = N_f \times T_{\text{frag}} = 9 \times 1{,}692 \approx 14.9~\text{s}
  \label{eq:Timage}
\end{equation}
matching the hardware measurement in Table~\ref{tab:e19_results}.

\subsubsection{Security Claim and Honest Framing}

We claim that \texttt{ENC\_DC\_ONLY} achieves \emph{visual-domain
confidentiality} under the following attacker model: the adversary intercepts
all LoRa packets, reassembles the fragment stream, and attempts JPEG decoding.
The stride corrupts the SOI marker and Huffman tables, preventing decoding.
The claim is \emph{not} that the scheme achieves IND-CPA security or hides
frequency-domain coefficients selectively --- it does not.  A determined
cryptanalyst with knowledge of the stride pattern and access to many
plaintext-ciphertext pairs could recover the keystream (since it is static
per-stride for a fixed image and frag counter), especially if the plaintext
header bytes are known.  The scheme's purpose is operational security against
passive image reconstruction, not semantic security.

\subsection{Fragment-Chained HMAC (FC-HMAC)}
\label{sec:fc_hmac}

\subsubsection{Problem Statement and Threat Model}

A 2{,}046-byte JPEG image spans 9 LoRa fragments.  The threat model for the
fragmentation layer is a Dolev-Yao adversary \cite{tamarin_prover} who:
\begin{itemize}
  \item Can intercept any LoRa packet transmitted in range.
  \item Can delay, replay, or reorder intercepted packets in any session.
  \item Can attempt to inject forged fragments with a chosen payload.
  \item Does \emph{not} know the session key $K$ (pre-shared, not on the radio).
\end{itemize}

Per-packet independent HMAC prevents the attacker from forging
\emph{individual} fragments but provides no protection against:
\begin{itemize}
  \item \textbf{Replay attack:} fragments from session $s$ replayed in session $s'$;
    each fragment passes its HMAC check independently.
  \item \textbf{Ordering attack:} fragment $i$ replaces fragment $j$; both have
    valid MACs, and the receiver has no way to detect the swap.
  \item \textbf{Cross-splice:} fragments from two sessions spliced together;
    each fragment is individually valid.
\end{itemize}
FC-HMAC addresses all three by binding each fragment's tag to the complete
prior chain state.

\subsubsection{FC-HMAC Construction}

Let $K$ be a 16-byte pre-shared session key.  Let $N_f$ be the total number
of fragments for one image.  Define the \emph{chain root} as:
\begin{equation}
  h_0 = \mathrm{HMAC\text{-}SHA256}(K,\;
        \mathtt{0x00000000} \;\|\; N_f)[\text{:}8]
  \label{eq:chain_root}
\end{equation}
For each fragment $i = 1, \ldots, N_f$, compute the \emph{chain tag}:
\begin{align}
  \mathit{tag}_i &= \mathrm{HMAC\text{-}SHA256}\!\Bigl(K,\;
    \underbrace{\mathit{frag\_id}_i}_{\text{2B}} \;\|\;
    \underbrace{\mathit{flags}_i}_{\text{1B}} \;\|\;
    \underbrace{\mathit{payload}_i}_{\text{up to 233B}} \;\|\;
    \underbrace{\mathit{tag}_{i-1}}_{\text{8B}}
  \Bigr)[\text{:}8]
  \label{eq:chain_tag}
\end{align}
The receiver maintains a running chain state $\hat{h}$ initialised to $h_0$.
Upon receiving fragment $i$, it recomputes $\hat{t}_i$ from $\hat{h}$ and the
received fields, then checks $\hat{t}_i \stackrel{?}{=} \mathit{tag}_i$.  If
the check passes, $\hat{h} \leftarrow \mathit{tag}_i$.  If it fails, the chain
is broken at position $i$, and all subsequent fragments cannot be chain-verified
(though they may be buffered for potential retransmit-based recovery).

\textbf{Why chaining achieves replay resistance.}  An adversary who replays
fragments from session $s$ in session $s'$ must forge tags consistent with the
chain state of session $s'$.  The chain state depends on $h_0^{(s')}$, which
depends on session $s'$'s unique context (when \texttt{img\_seq} is included
in the root --- see Section~\ref{sec:limit_f3}).  Without knowing $K$, the
adversary cannot compute consistent $\mathit{tag}_i^{(s')}$ values.

\textbf{Why chaining achieves ordering resistance.}  Swapping fragment $j$
into position $i \neq j$ changes \texttt{frag\_id} and inserts a different
payload into the HMAC input at position $i$.  The recomputed tag will not
match the transmitted $\mathit{tag}_i$, which was computed with the correct
\texttt{frag\_id} $= i$ and its own payload.

\subsubsection{Chain Initialisation and img\_seq}

Equation~(\ref{eq:chain_root}) uses a zero nonce, which is the implementation
in the E19 firmware.  A production-ready implementation should include a
session-unique image sequence number \texttt{img\_seq} to bind the chain
root to a specific session:
\begin{equation}
  h_0^{(\text{prod})} = \mathrm{HMAC\text{-}SHA256}(K,\;
        \mathtt{img\_seq} \;\|\; N_f)[\text{:}8]
  \label{eq:chain_root_prod}
\end{equation}
The Tamarin formal model (Section~\ref{sec:tamarin}) uses
Equation~(\ref{eq:chain_root_prod}), and the F3 (No\_Replay) lemma is proved
under this stronger construction.  The E19 firmware gap (omitting
\texttt{img\_seq}) is documented as a known limitation in
Section~\ref{sec:limit_f3}.

\subsubsection{Security Parameters and Truncation Analysis}

The HMAC output is truncated to 8 bytes (64 bits).  Table~\ref{tab:hmac_trunc}
analyses the security-overhead tradeoff for three candidate truncation lengths:

\begin{table}[ht]
\centering
\caption{HMAC Truncation Length: Security vs.\ Overhead Tradeoff}
\label{tab:hmac_trunc}
\begin{tabular}{lrrl}
\toprule
Tag size & Overhead & Forgery prob. & Assessment \\
\midrule
4\,B (32-bit) & 1.7\% & $2^{-32}$ & Insufficient \\
8\,B (64-bit) & 3.3\% & $2^{-64}$ & \textbf{Target (selected)} \\
12\,B (96-bit) & 5.0\% & $2^{-96}$ & Strong but wasteful \\
\bottomrule
\end{tabular}
\end{table}

Under LoRa duty-cycle constraints, a single adversarial node can inject at
most $\approx$100 packets per hour.  At this injection rate, exhaustive search
against a 64-bit tag would require $\approx 1.8 \times 10^{17}$~years in
expectation --- well beyond any practical threat.

\subsection{Packet Format and Frame Layout}
\label{sec:pkt_format}

Each LoRa packet transmitted by the FC-HMAC protocol carries the following
fixed-layout frame:

\begin{table}[ht]
\centering
\caption{FC-HMAC Packet Frame Layout}
\label{tab:pkt_format}
\begin{tabular}{lrp{4.5cm}}
\toprule
Field & Width & Description \\
\midrule
\texttt{magic}        & 2\,B & Framing bytes: 0xE1 0x9A \\
\texttt{fragment\_id} & 2\,B & Zero-indexed fragment number (0 to $N_f{-}1$) \\
\texttt{total\_frags} & 2\,B & Total fragments in this image \\
\texttt{flags}        & 1\,B & Bit 0: last fragment; bits 1--2: enc mode \\
\texttt{payload}      & $N$\,B & Possibly encrypted image bytes \\
\texttt{hmac8}        & 8\,B & 8-byte truncated chained HMAC tag \\
\midrule
\textbf{Max total}    & \multicolumn{2}{l}{$7 + N + 8 \leq 255$\,B; $N \leq 240$\,B} \\
\textbf{Used ($N=233$)} & \multicolumn{2}{l}{248\,B total per packet} \\
\bottomrule
\end{tabular}
\end{table}

The 7-byte fixed header plus 8-byte HMAC tag consume 15~bytes, leaving 240~B
for payload.  In practice the E19 firmware uses $N = 233$~bytes payload for
a 248-byte total, providing a 7-byte margin below the 255-byte hardware limit.
The 8-byte \texttt{hmac8} field represents 3.3\% overhead relative to the
241-byte application-layer frame (7 header + 234~B payload before HMAC).

The \texttt{magic} bytes serve as a lightweight frame synchronisation marker:
the receiver discards any packet that does not begin with 0xE1~0x9A, preventing
stray LoRa broadcasts from polluting the fragment reassembly buffer.

\subsubsection{Fragment Reassembly at the Receiver}

The receiver maintains the following state per ongoing image transmission:
\begin{itemize}
  \item \texttt{total\_frags}: total fragment count, learned from the first
    received fragment's \texttt{total\_frags} field.
  \item \texttt{recv\_bitmap}: 64-bit bitmask (supports up to 64 fragments)
    indicating which \texttt{fragment\_id} values have been received.
  \item \texttt{chain\_state} ($\hat{h}$): the 8-byte chain tag of the last
    successfully chain-verified fragment.  Initialised to $h_0$
    (Equation~\ref{eq:chain_root}) before the first fragment.
  \item \texttt{reassembly\_buffer}: array of $N_f$ payload slots, each 233~B.
  \item \texttt{chain\_broken}: boolean flag set when a chain verification
    failure is detected; once set, subsequent fragments cannot be
    chain-verified but are still stored for diagnostic purposes.
\end{itemize}

Algorithm~\ref{alg:receiver} gives the per-fragment receive procedure:

\begin{algorithm}[H]
\caption{FC-HMAC Fragment Reception and Chain Verification}
\label{alg:receiver}
\begin{algorithmic}[1]
\Require Received packet bytes \textit{pkt}
\State Parse \texttt{magic}, \texttt{frag\_id}, \texttt{total\_frags},
       \texttt{flags}, \texttt{payload}, \texttt{hmac8}
\If{\texttt{magic} $\neq$ \texttt{0xE19A}} \Return DISCARD \EndIf
\If{\texttt{frag\_id} $\geq$ \texttt{total\_frags}} \Return DISCARD \EndIf
\State $\hat{t} \leftarrow \mathrm{HMAC\text{-}SHA256}(K,\;
       \langle \texttt{frag\_id}, \texttt{flags}, \texttt{payload}, \hat{h} \rangle)[\text{:}8]$
\If{$\hat{t} \neq \texttt{hmac8}$}
  \State Set \texttt{chain\_broken} $\leftarrow$ \texttt{true}
  \State Record chain-break at \texttt{frag\_id}; \Return CHAIN\_BREAK
\EndIf
\State $\hat{h} \leftarrow \texttt{hmac8}$ \Comment{Advance chain state}
\State Copy \texttt{payload} to \texttt{reassembly\_buffer[frag\_id]}
\State Set bit \texttt{frag\_id} in \texttt{recv\_bitmap}
\If{\texttt{recv\_bitmap} $=$ $(1 \ll \texttt{total\_frags}) - 1$}
  \State Deliver assembled image to application \Return COMPLETE
\EndIf
\Return OK
\end{algorithmic}
\end{algorithm}

When all bits in \texttt{recv\_bitmap} are set, the image is declared complete
and passed to the application layer.  A missing fragment (due to channel loss)
breaks the chain from that position onward; the receiver can request a
retransmit using the bitmap information (in an ACK-capable variant) or accept
a partial image by re-anchoring the chain at the next received fragment.

The \texttt{chain\_broken} flag is important for diagnostics: it allows the
receiver to distinguish between (a)~a complete image received without chain
breaks, (b)~a complete image received with at least one chain break (indicating
tampering or a firmware bug), and (c)~an incomplete image (missing fragments).
In the E19 experiment, case~(a) was observed for all 8 transmitted images.

%% -----------------------------------------------------------------------
\section{Security Analysis}
\label{sec:security}

\subsection{Tamarin Formal Model}
\label{sec:tamarin}

\subsubsection{Model Overview}

We model FC-HMAC in Tamarin Prover \cite{tamarin_prover} (version 1.13.0,
executed in the \texttt{lora-mesh-tamarin} Docker container).  The model file
is \texttt{phase2\_tamarin\_models/paper\_d/fc\_hmac.spthy} (verification log:
\texttt{phase2\_tamarin\_models/paper\_d/fc\_hmac\_verification.log}).  The
model encodes a two-fragment session to keep the state space finite while
capturing all relevant security properties.

\subsubsection{Modelling Decisions}

\textbf{MAC as perfect PRF.}  The \texttt{mac/2} function symbol is declared
with no inversion or cancellation equations.  This models HMAC-SHA256 as an
ideal pseudorandom function under the standard symbolic (Dolev-Yao) model,
where the adversary cannot invert, decompose, or partially evaluate the MAC
without the key.

\textbf{Dolev-Yao network adversary.}  The attacker controls all network
communication: it can intercept, store, replay, and inject arbitrary LoRa
messages.  It knows all public values (fragment IDs, total-frag counts,
ciphertexts) but not the session key $K$.

\textbf{State-space bounds.}  Two Tamarin restrictions limit rule applications:
\texttt{OneSetup} ensures a single sender initialisation per session, and
\texttt{ReceiverOnce} bounds receiver state to one verification chain per
session.  These restrictions prevent unbounded rule application while preserving
the security properties of interest.

\textbf{Chain initialisation.}  The \texttt{Sender\_Init} rule computes
$h_0 = \mathit{mac}(K,\, \langle \mathit{sess},\, \mathit{tf} \rangle)$,
where \texttt{sess} is a fresh session identifier and \texttt{tf} is the total
fragment count.  This matches Equation~(\ref{eq:chain_root_prod}).

\textbf{Fragment rules.}  \texttt{Sender\_Frag1} and \texttt{Sender\_Frag2}
each compute a chained tag
$\mathit{tag}_i = \mathit{mac}(K,\,
 \langle \mathit{frag\_id},\, f,\, p,\, \mathit{tag}_{\mathit{prev}} \rangle)$
and output the packet to the adversary-controlled network via \texttt{Out(\ldots)}.
\texttt{Receiver\_Frag1} and \texttt{Receiver\_Frag2} each receive a packet
via \texttt{In(\ldots)}, re-derive the expected tag, and enforce equality via
the \texttt{Eq(lhs, rhs)} restriction, which causes Tamarin to prune all
execution traces where the check fails.

\subsubsection{Lemma Statements and Results}

The four lemmas and their proof results are summarised in
Table~\ref{tab:tamarin}.  All four lemmas were proved in 1.53~seconds total on a
standard workstation (Intel Core i7, 16~GB RAM) running Tamarin 1.13.0.

\begin{table}[ht]
\centering
\caption{Tamarin Verification Results (fc\_hmac.spthy, 1.53\,s total)}
\label{tab:tamarin}
\begin{tabular}{cllr}
\toprule
ID & Lemma Name & Result & Steps \\
\midrule
F4 & \textit{Executability} & PROVED & 14 \\
F1 & \textit{Fragment\_Authenticity} & PROVED & 10 \\
F2 & \textit{Chain\_Gap\_Detectable} & PROVED & 3 \\
F3 & \textit{No\_Replay} & PROVED & 12 \\
\bottomrule
\end{tabular}
\end{table}

\textbf{F4 (Executability, 14 steps)} is a sanity lemma: it proves the
existence of a complete honest execution trace where the sender sends two
fragments and the receiver successfully verifies both.  This confirms the
model is not vacuously satisfied.

\textbf{F1 (Fragment\_Authenticity, 10 steps)} states: if the receiver
successfully verifies fragment $i$, then there exists a sender action that
emitted a fragment with the same fragment ID, flags, payload, and chain tag.
Formally, the adversary cannot cause the receiver to accept a fragment that
was never sent by a legitimate sender holding key $K$.

\textbf{F2 (Chain\_Gap\_Detectable, 3 steps)} states: if fragment $i$ is
missing from the chain and fragment $i+1$ is received, then the receiver's
verification of fragment $i+1$ will fail.  The chain cannot ``skip'' a missing
fragment without detection.  This property is unique to FC-HMAC and has no
analogue in per-packet independent HMAC schemes.

\textbf{F3 (No\_Replay, 12 steps)} states: if the receiver accepts fragment
$i$ from session $s'$, then it was originally sent in session $s'$ (not
replayed from a different session $s \neq s'$).  This property requires
\texttt{img\_seq} in the chain root, as discussed in
Section~\ref{sec:limit_f3}.

\subsubsection{Security Properties Summary}

Table~\ref{tab:security_props} summarises the four security properties in the
FC-HMAC model and maps them to the design elements and attacker goals they address.

\begin{table}[ht]
\centering
\caption{FC-HMAC Security Properties and Design Mapping}
\label{tab:security_props}
\begin{tabular}{cllp{2.4cm}}
\toprule
Lemma & Property & Precondition & Defeated attack \\
\midrule
F1 & Fragment authenticity & Key $K$ unknown to adversary & Fragment forgery \\
F2 & Chain-gap detectable & Chain integrity & Silent fragment drop \\
F3 & No cross-session replay & \texttt{img\_seq} in chain root & Session replay \\
F4 & Executability (witness) & None & Protocol liveness \\
\bottomrule
\end{tabular}
\end{table}

\subsubsection{Comparison with Independent Per-Packet HMAC}

To underscore the novelty of chain-based verification, Table~\ref{tab:hmac_compare}
contrasts the security properties achievable by independent per-packet HMAC
(as in Ryczek \cite{ryczek2025}) versus FC-HMAC.

\begin{table}[ht]
\centering
\caption{Independent Per-Packet HMAC vs.\ FC-HMAC: Security Property Coverage}
\label{tab:hmac_compare}
\begin{tabular}{lcc}
\toprule
Security Property & Per-Packet HMAC & FC-HMAC \\
\midrule
Fragment forgery prevention     & \checkmark & \checkmark \\
Individual packet integrity     & \checkmark & \checkmark \\
Fragment reorder detection      & $\times$   & \checkmark \\
Fragment insertion detection    & $\times$   & \checkmark \\
Cross-session replay prevention & $\times$   & \checkmark (with \texttt{img\_seq}) \\
Silent gap detection            & $\times$   & \checkmark \\
Formally verified (machine-checked) & $\times$ & \checkmark \\
\bottomrule
\end{tabular}
\end{table}

\subsubsection{F3 Firmware Gap}
\label{sec:limit_f3_tamarin}

The Tamarin F3 lemma is proved under the model with \texttt{img\_seq} in the
chain root.  The E19 firmware (Section~\ref{sec:e19}) uses a zero nonce
instead, creating a gap between the formal model and the hardware
implementation.  This is a documented limitation; the E20 firmware (planned)
addresses it by encoding \texttt{img\_seq} in bytes 0--3 of the
\texttt{frag\_id=0} payload prefix.  The F1, F2, and F4 lemmas hold
regardless of whether \texttt{img\_seq} is present.

\subsection{Visual Security Evaluation (E19-PSNR)}
\label{sec:psnr}

\subsubsection{Evaluation Methodology}

The visual security evaluation script (\texttt{tools/psnr\_attack\_eval.py})
applies each encryption mode to a synthetic JPEG and then simulates attacker
reconstruction attempts.  Two attacker models are considered:

\textbf{A1 --- Naive attacker:} receives the encrypted byte stream
(after FC-HMAC verification strips the header) and attempts to decode it
directly as a JPEG file using a standard decoder.

\textbf{A2 --- Format-aware attacker:} knows the stride pattern (16 bytes
per 64-byte window) and the starting offset.  Before decoding, the attacker
zeroes all bytes at stride positions, replacing ciphertext with zeros to
recover the underlying format structures.  This simulates a partial chosen-plaintext
attack by an adversary with full knowledge of the encryption algorithm but
not the key.

For both attackers, success is measured as either a finite PSNR value (image
decoded, compared to original) or a JPEG decode failure.

\subsubsection{Results and Analysis}

\begin{table}[ht]
\centering
\caption{Visual Security Evaluation Results (E19-PSNR, \texttt{psnr\_attack\_eval.py})}
\label{tab:psnr}
\begin{tabular}{cllll}
\toprule
Mode & Enc bytes & Stride calls & A1 & A2 \\
\midrule
\texttt{ENC\_NONE}    & 0\%   & 0 & PSNR $= \infty$ & N/A \\
\texttt{ENC\_DC\_ONLY}& 24.9\% & 9 & \textbf{Decode FAIL} & \textbf{Decode FAIL} \\
\texttt{ENC\_FULL}    & 99.8\% & 256 & \textbf{Decode FAIL} & \textbf{Decode FAIL} \\
\bottomrule
\end{tabular}
\end{table}

\texttt{ENC\_DC\_ONLY} at 24.9\% encrypted bytes completely prevents JPEG
reconstruction under both attacker models.  The A2 format-aware attacker fails
because the stride begins at byte~0 (SOI marker) and the 9~encryption windows
cover the first $9 \times 64 = 576$~bytes of the image.  Within this range,
the JPEG SOI marker (bytes 0--1), APP0 header (bytes 2--20), quantisation
tables (bytes $\sim$21--200), and Huffman tables (bytes $\sim$200--500) are all
partially encrypted.  Even if the attacker zeroes the encrypted positions,
the non-zero residual pattern in the surrounding bytes is inconsistent with any
valid JPEG structure, causing decoder failure.

The PSNR target originally set for this evaluation (PSNR attack $<$ 20~dB) is
superseded: decode failure provides a stronger guarantee than a finite PSNR
value, since the attacker obtains zero visual information rather than a
degraded reconstruction.

\subsubsection{Limitation: Semantic Security}

The E19-PSNR evaluation confirms that format-level obfuscation is achieved
against both attacker models tested.  However, we do not claim semantic
security (IND-CPA) for the stride scheme.  An adversary with access to multiple
plaintexts encrypted under the same key and stride would observe the same XOR
pattern at stride positions across images, potentially recovering keystream
blocks.  The scheme is intended for deployments where the same key is rotated
periodically and the attacker has limited oracle access.

%% -----------------------------------------------------------------------
\section{Hardware Evaluation}
\label{sec:eval}

\subsection{Platform and Experimental Setup}
\label{sec:platform}

Both transmitter and receiver nodes use the EoRa PI development board, which
integrates an Espressif ESP32-S3 \cite{esp32s3_datasheet} (Xtensa LX7 dual-core,
240~MHz, 512~KB internal SRAM, 2~MB PSRAM) with a Semtech SX1262 LoRa
transceiver \cite{semtech_sx1262} on-board.  The SX1262 is connected via SPI
with the following pin mapping: NSS = GPIO7, SCK = GPIO5, MOSI = GPIO6,
MISO = GPIO3, RST = GPIO8, BUSY = GPIO34, DIO1 = GPIO33.

The firmware is compiled with ESP32 Arduino SDK~3.x with CDCOnBoot enabled for
USB serial output.  The mbedTLS library bundled with the SDK provides
software-only HMAC-SHA256 and AES operations; no hardware cryptographic
accelerator is used.  All CPU timing measurements reflect pure software
execution on a single Xtensa LX7 core.

Experimental parameters are summarised in Table~\ref{tab:platform}.

\begin{table}[ht]
\centering
\caption{Hardware Platform and LoRa Configuration (E19)}
\label{tab:platform}
\begin{tabular}{ll}
\toprule
Parameter & Value \\
\midrule
MCU & ESP32-S3 (Xtensa LX7, 240\,MHz, 2\,MB PSRAM) \\
Radio & Semtech SX1262 \\
Firmware SDK & ESP32 Arduino SDK 3.x, mbedTLS \\
Frequency & 922.0\,MHz \\
Spreading Factor & SF9 \\
Bandwidth & 125\,kHz \\
Coding Rate & 4/7 \\
TX power & 2\,dBm \\
Max LoRa payload & 255\,B (SX1262 hardware register limit) \\
Fragment size & 248\,B (7\,B header + 233\,B data + 8\,B HMAC) \\
Session key & 16-byte pre-shared (AES-128 equivalent) \\
Test image & Synthetic JPEG, $200\times150$ px, 2{,}046\,B \\
Fragments per image & 9 \\
Measured fragment airtime & 1{,}692\,ms \\
Measured image airtime & 14.9\,s \\
Experiment duration & 120\,s \\
Images transmitted & 8 (complete) \\
\bottomrule
\end{tabular}
\end{table}

The transmitter sends the test image as a C byte array (\texttt{test\_image.h})
embedded in the firmware flash.  This approach isolates the protocol layer from
camera hardware variability, enabling reproducible measurements.  A live camera
attachment (OV2640, DVP interface via EoRa PI GPIO) is planned for Phase~3
validation (Section~\ref{sec:future}).

\subsection{E19 Experiment Results}
\label{sec:e19}

Table~\ref{tab:e19_results} reports all primary measurements from experiment
E19 (results file: \texttt{hardware\_experiments/E19\_selective\_enc/results/e19\_results.md}).

\begin{table}[ht]
\centering
\caption{E19 Hardware Experiment Summary (120\,s, 2$\times$ EoRa PI, SF9/BW125)}
\label{tab:e19_results}
\begin{tabular}{ll}
\toprule
Metric & Measured Value \\
\midrule
Images transmitted & 8 \\
Total fragments & 72 \\
Fragments verified (chain\_ok) & 72 / 72 = \textbf{100\%} \\
Packet Delivery Ratio (PDR) & \textbf{100\%} \\
Airtime per image & 14.9\,s \\
Measured fragment airtime & 1{,}692\,ms per fragment \\
\midrule
HMAC overhead (per packet) & 8\,B / 241\,B = \textbf{3.3\%} \\
HMAC overhead (per image) & 72\,B / 2{,}046\,B = \textbf{3.5\%} \\
HMAC CPU per image & $\approx$950\,$\mu$s (9 HMAC-SHA256 calls) \\
Enc CPU per image (DC\_ONLY) & 2{,}807\,$\mu$s (9 HMAC-PRF keystream calls) \\
Total crypto CPU per image & 950 + 2{,}807 = \textbf{3{,}757\,$\mu$s} ($\approx$ 3.76\,ms) \\
\midrule
Encrypted bytes & 496 / 2{,}046\,B = \textbf{24.2\%} \\
Encryption mode & \texttt{ENC\_DC\_ONLY} \\
\bottomrule
\end{tabular}
\end{table}

All 72 fragments in the 120-second run pass FC-HMAC chain verification under
benign conditions (no packet loss, no tampering, line-of-sight 1-metre indoor
range).  The 100\% chain-ok rate confirms that (i)~the HMAC-SHA256 computation
is deterministic and consistent between sender and receiver, (ii)~the chain
initialisation and tag update logic are correctly implemented, and (iii)~the
SX1262 LoRa radio delivers fragments without bit errors in this test
environment.

The HMAC overhead of 3.3\% per packet is below the theoretical 4-byte-tag
bound of 1.7\%, which would be insufficient for security, and below the
12-byte bound of 5.0\%, which would be wasteful --- placing the 8-byte design
squarely at the optimal security-overhead point (Table~\ref{tab:hmac_trunc}).

\subsubsection{Fragment Timing and Inter-Fragment Gap}

Table~\ref{tab:frag_timing} breaks down the measured per-fragment timing into
its constituent phases as observed in the E19 firmware.

\begin{table}[ht]
\centering
\caption{Per-Fragment Timing Breakdown (E19, 248-byte fragment, SF9/BW125)}
\label{tab:frag_timing}
\begin{tabular}{lr}
\toprule
Phase & Duration \\
\midrule
SX1262 TX preparation (SPI config) & $\approx$5\,ms \\
LoRa preamble + header on air & 50.9\,ms \\
Payload on air (107 symbols $\times$ 4.096\,ms) & 438\,ms \\
TX-to-idle turnaround & $\approx$2\,ms \\
\midrule
\textit{Theoretical packet on air} & \textit{489\,ms} \\
\midrule
Application-layer inter-fragment gap & $\approx$1{,}200\,ms \\
\textbf{Total per-fragment cycle (measured)} & \textbf{1{,}692\,ms} \\
\bottomrule
\end{tabular}
\end{table}

The 1{,}200~ms inter-fragment gap is enforced at the application layer to
comply with the ETSI 1\% duty-cycle rule: $T_{\text{on-air}} /
(T_{\text{on-air}} + T_{\text{gap}}) = 489 / (489 + 1{,}200) \approx 0.29\%$,
well below the 1\% ceiling.  The gap also gives the receiver time to process
the received fragment and update its chain state before the next packet arrives.

Within the 1{,}200-ms inter-fragment gap, the sender computes the selective
encryption keystream and the FC-HMAC tag for the next fragment (3{,}757~$\mu$s /
9 = 417~$\mu$s per fragment), builds the packet frame, and loads it into the
SX1262 FIFO.  The crypto processing occupies only 0.035\% of the inter-fragment
gap, confirming that FC-HMAC introduces no detectable increase in duty-cycle
compliance risk.

\subsection{Efficiency Analysis}
\label{sec:efficiency}

\subsubsection{Airtime Invariance Across Encryption Modes}

A fundamental observation from Table~\ref{tab:enc_modes} is that selective
encryption does \emph{not} reduce airtime relative to full encryption.  All
three modes (\texttt{ENC\_NONE}, \texttt{ENC\_DC\_ONLY}, \texttt{ENC\_FULL})
transmit 9 fragments of 248~bytes each, producing identical 14.9-second airtime.

This invariance arises because airtime is determined by:
\begin{equation}
  T_{\text{image}} = N_f \times T_{\text{frag}}(L_{\text{frag}})
  \label{eq:airtime}
\end{equation}
where $N_f = \lceil N_{\text{img}} / N_{\text{payload}} \rceil$ depends only
on image size and per-fragment payload capacity, and $T_{\text{frag}}$ depends
only on fragment length and LoRa parameters.  Neither depends on how many bytes
within the payload are encrypted.  Encryption changes individual byte values
but not the byte count.

The practical implication for the design: selective encryption is \emph{not}
an airtime optimisation technique.  It is a CPU optimisation technique.
Claims in the literature that selective encryption reduces LoRa duty-cycle
consumption should be scrutinised to determine whether they refer to CPU
cost or to selective \emph{fragment omission} (transmitting fewer fragments
by skipping non-critical image regions).

\subsubsection{CPU Cost Breakdown}

\begin{table}[ht]
\centering
\caption{CPU Cost by Encryption Mode (2{,}046-byte image, ESP32-S3 240\,MHz)}
\label{tab:cpu_savings}
\begin{tabular}{lrrrr}
\toprule
Mode & Enc CPU & HMAC CPU & Total crypto & vs.\ FULL \\
\midrule
\texttt{ENC\_NONE}    & 0\,$\mu$s & 950\,$\mu$s & 950\,$\mu$s & $-$92\% \\
\texttt{ENC\_DC\_ONLY}& 2{,}807\,$\mu$s & 950\,$\mu$s & 3{,}757\,$\mu$s & $-$\textbf{76\%} \\
\texttt{ENC\_FULL}    & $\approx$11{,}600\,$\mu$s & 950\,$\mu$s & $\approx$12{,}550\,$\mu$s & baseline \\
\bottomrule
\end{tabular}
\end{table}

The \texttt{ENC\_DC\_ONLY} mode reduces encryption CPU by 76\% compared to
\texttt{ENC\_FULL} (2{,}807~$\mu$s vs.\ $\approx$11{,}600~$\mu$s) while
maintaining format-level visual obfuscation (Table~\ref{tab:psnr}).  The HMAC
CPU (950~$\mu$s) is identical across all modes because all fragments carry the
FC-HMAC chain tag.

Total crypto CPU for \texttt{ENC\_DC\_ONLY} is 3{,}757~$\mu$s (3.76~ms) per
2{,}046-byte image.  For a one-image-per-minute transmission schedule, the image
transmission occupies 14.9~s of airtime and 3.76~ms of crypto CPU per 60-second
cycle.  The crypto fraction of the non-airtime processing budget is therefore
negligible: $3.76 \text{ ms} / (60\,\text{s} - 14.9\,\text{s}) \approx 0.008\%$.

\subsubsection{Overhead Scalability}

The FC-HMAC overhead has two components:
\begin{enumerate}
  \item \textbf{Per-packet tag overhead:} 8~B per fragment, constant regardless
    of payload size.  For a 233-byte payload, this is 3.3\%.  For a
    30-byte payload (in a network where packets are smaller), this rises to
    21\% --- a significant penalty.  The 8-byte tag is therefore best suited
    to near-maximum-payload fragments.
  \item \textbf{Per-image CPU overhead:} $O(N_f)$ HMAC calls, one per fragment.
    For a 20-kB image (as in Ryczek \cite{ryczek2025}) split into
    $\lceil 20480 / 233 \rceil = 88$~fragments, the HMAC CPU scales to
    $\approx$9{,}250~$\mu$s --- still below the estimated 113{,}600~$\mu$s for
    full encryption of 20~kB, representing the same 76\% relative saving.
\end{enumerate}

%% -----------------------------------------------------------------------
\section{Related Work}
\label{sec:related}

\subsection{Image Transmission Security over LoRa}

Ryczek and Sobieraj \cite{ryczek2025} present the most closely related prior
work: an ESP32-CAM paired with an E220 LoRa module (SX1262-based) transmitting
20~kB camera images with full AES-128-CBC encryption and independent
per-packet HMAC-SHA256.  They report PDR~=~100\% in indoor experiments and
$<$1\% energy overhead on the transmitter side for encrypted vs.\ plaintext
transmission.  Their work demonstrates the viability of application-layer
encryption on ESP32-class hardware.

Our work differs in three structural ways.  First, we use \emph{selective}
rather than full encryption, demonstrating a 76\% CPU saving while maintaining
visual obfuscation.  Second, we replace independent per-packet HMAC with a
\emph{chained} structure (FC-HMAC) that binds the entire fragment sequence,
enabling detection of ordering attacks and cross-session replay.  Third, we
provide machine-checked formal proofs of chain-gap detectability and no-replay
via Tamarin Prover --- properties that per-packet HMAC schemes cannot satisfy
and that have not been formally verified in prior LoRa image security work.

\subsection{Selective Encryption for Compressed Media}

Selective encryption has been extensively studied for MPEG video
\cite{rfc3610}: encrypting I-frame headers, motion vectors, or DC coefficients
provides visual confidentiality at a fraction of full-encryption cost.  The
key insight is that video and image codecs have an \emph{entropy hierarchy}:
a small number of structurally critical bits, when corrupted, propagate errors
throughout the decoded signal.  Targeting these bits maximises disruption per
encrypted byte.

Applying this principle to JPEG over LoRa introduces a new constraint absent
from video settings: the encoder runs on a microcontroller without floating-point
DSP hardware, the JPEG bitstream is static (not a live stream), and the
per-fragment overhead budget is measured in bytes, not kilobytes.  Our
stride-based approximation is a pragmatic adaptation that achieves the same
format-level disruption without requiring a Huffman-stream parser.

Earlier work on selective JPEG encryption for mobile and embedded devices
targeted megapixel-range images on processors with hardware floating-point and
large caches, where full JPEG parsing is feasible.  At the LoRa constraint
level (2~KB images, 240-MHz MCU with no hardware AES, 8-byte overhead budget per
packet), a parsing-free approximation is the appropriate design point.

It is important to note that selective encryption of JPEG images can achieve
format disruption at the \emph{syntax} layer (as in our scheme) or at the
\emph{semantic} layer (targeting specific DCT frequencies).  Our scheme operates
at the syntax layer: by corrupting the Huffman table segments, we prevent the
decoder from interpreting any subsequent coded data.  This is a weaker claim
than frequency-domain selective encryption but is achievable with a 9-call
keystream computation in firmware.  The relationship between syntax-layer and
frequency-domain selective encryption is discussed further in
Section~\ref{sec:limit_b1}.

\subsection{MAC Authentication for Constrained Wireless Networks}

TinySec \cite{tinysec} introduced MAC-only and encryption+MAC modes for IEEE
802.15.4 sensor networks in 2004, demonstrating that 4-byte MACs are sufficient
for low-data-rate wireless links at acceptable energy overhead.  TinySec uses
CBC-MAC (cipher-block chaining MAC) applied independently to each packet,
which prevents forgery of individual packets but does not chain across packet
boundaries.  TinySec's design inspired the choice of compact MAC sizes in
constrained networks; our 8-byte tag selection extends this by adding
cross-packet chain binding.

RFC~3610 \cite{rfc3610} defines AES-CCM, which combines encryption and
authentication in a single pass using CTR mode for encryption and CBC-MAC for
authentication.  AES-CCM requires hardware AES support (or significant software
overhead) and does not natively support fragment chaining.  Furthermore,
AES-CCM is not available in mbedTLS without an AES hardware accelerator on
ESP32-S3 (the hardware AES driver is not linked in the ESP32 Arduino SDK
configuration used in E19).  Our scheme uses HMAC-SHA256 (available as a
pure-software primitive in mbedTLS), truncated to 8 bytes, chained across
fragments.

The HMAC-based PRF construction \cite{bellare_hmac} provides a provable
security reduction from the MAC security to the PRF security of the underlying
hash function (SHA-256), under the random oracle model.  This is a stronger
foundation than CBC-MAC (which has known weaknesses for variable-length
messages) and more widely available in embedded TLS stacks than AES-CMAC.

\subsection{Formal Verification of IoT Protocol Security}

Tamarin Prover \cite{tamarin_prover} has been applied to verify LoRaWAN
join-accept procedures \cite{tamarin_prover}, lightweight RFID authentication
protocols, and TLS 1.3.  In the LoRa context, prior Tamarin work has focused
on the LoRaWAN session-establishment phase, not on application-layer
fragmentation integrity.  Our FC-HMAC model is the first formal verification
of a fragment-chained MAC for image transmission over LoRa.

The F2 (Chain\_Gap\_Detectable) and F3 (No\_Replay) lemmas are novel
contributions: F2 captures a property that is unique to chained MACs and has
no analogue in per-packet HMAC systems; F3 proves session binding that is
typically assumed but not formally checked in IoT security papers.

\subsection{Detailed Comparison: FC-HMAC vs.\ Ryczek \& Sobieraj 2025}

Table~\ref{tab:comparison} provides a structured seven-dimension comparison
between FC-HMAC (this work) and the nearest prior work \cite{ryczek2025}.

\begin{table*}[ht]
\centering
\caption{FC-HMAC vs.\ Ryczek \& Sobieraj 2025: Seven-Dimension Comparison}
\label{tab:comparison}
\begin{tabular}{p{2.6cm}p{5.0cm}p{5.8cm}}
\toprule
Dimension &
  Ryczek \& Sobieraj (Electronics 2025) \cite{ryczek2025} &
  \textbf{FC-HMAC (this work)} \\
\midrule
\textbf{Encryption scope} &
  Full image, AES-128-CBC, 100\% bytes encrypted &
  Stride-based selective (24.2\% bytes); format-level obfuscation; 76\% CPU saving \\
\addlinespace
\textbf{Authentication} &
  Per-packet HMAC-SHA256, independent tags; each packet self-contained &
  FC-HMAC chain: $\mathit{tag}_i$ commits to $\mathit{tag}_{i-1}$; reordering, insertion, replay detectable \\
\addlinespace
\textbf{Loss behaviour} &
  Lost packet does not affect other packets; receiver can accept partial image transparently &
  Chain break at loss position; receiver obtains a verified prefix chain up to the first gap \\
\addlinespace
\textbf{Formal verification} &
  None &
  Tamarin Prover: F1 (authenticity), F2 (chain-gap), F3 (no-replay), F4 (executability); all proved in 1.53\,s \\
\addlinespace
\textbf{HMAC overhead} &
  $\approx$4--5\% (estimated from paper data; not directly hardware-measured) &
  3.3\% per packet, 3.5\% per image; directly hardware-measured in E19 \\
\addlinespace
\textbf{PSNR under attack} &
  Not measured; visual security assumed from full encryption &
  JPEG decode FAIL for both A1 (naive) and A2 (format-aware) attackers; measured via \texttt{psnr\_attack\_eval.py} \\
\addlinespace
\textbf{Hardware platform} &
  ESP32-CAM + EBYTE E220 (SX1262), 20\,kB JPEG images, OV2640 camera &
  EoRa PI (ESP32-S3 + SX1262), 2{,}046-byte synthetic JPEG; OV2640 attachment planned for Phase~3 \\
\bottomrule
\end{tabular}
\end{table*}

%% -----------------------------------------------------------------------
\section{Discussion}
\label{sec:discussion}

\subsection{Limitation B1: Stride Approximation vs.\ True DCT-Domain Encryption}
\label{sec:limit_b1}

The primary theoretical limitation of Phase~1 is that the stride-based scheme
is \emph{not} true frequency-domain selective encryption.  We explicitly
acknowledge and reframe this as follows.

True selective encryption for JPEG --- specifically targeting DCT DC coefficient
bytes --- requires streaming through the Huffman-coded bitstream, which encodes
DC differences as variable-length codes that do not occur at predictable byte
offsets.  Implementing such a parser on a constrained embedded node introduces
CPU overhead comparable to full-image encryption for the decoding step alone.

The stride-based approximation achieves visual-domain disruption through a
different mechanism: it corrupts the JPEG format infrastructure (SOI marker,
Huffman tables) rather than DCT frequency coefficients.  The E19-PSNR evaluation
confirms that this achieves the operational goal (JPEG decode failure under both
attacker models) without frequency-domain targeting.  The honest framing of the
contribution is therefore:

\textit{Lightweight visual obfuscation via structure-level disruption, not
frequency-domain selective encryption.}

Future work (\texttt{tools/jpeg\_dc\_offset\_extractor.py}) will quantify the
actual overlap between stride positions and true DC coefficient byte positions
for a corpus of JPEG images.  If the overlap is high ($>$50\%), the
approximation claim is partially vindicated; if it is low ($<$20\%), the
obfuscation result must be attributed entirely to header corruption, and the
``DCT-targeting'' framing should be abandoned.

The path to true frequency-domain selective encryption for Phase~2 would
require: (a)~a lightweight JPEG entropy decoder that identifies DC code
boundaries, (b)~in-place XOR of those bytes only, and (c)~evaluation of
whether the decoder overhead is justified by the additional security properties
(IND-CPA against an attacker who knows the stride pattern but not the key).

\subsection{Limitation F3: img\_seq Omission in E19 Firmware}
\label{sec:limit_f3}

The Tamarin F3 lemma (No\_Replay) is proved under the assumption that
\texttt{img\_seq} is included in the chain root (Equation~\ref{eq:chain_root_prod}).
The E19 firmware uses a zero nonce (Equation~\ref{eq:chain_root}), creating a
gap between the formal model and the hardware implementation.

Under the E19 firmware, an adversary who captures a complete valid fragment
sequence from session $s$ and replays it as session $s'$ will succeed: the
receiver re-derives the same chain tags (since $h_0$ is the same for both
sessions) and accepts the replayed fragments.  This attack requires the
adversary to capture a complete prior session in its entirety, which is a
non-trivial requirement in a LoRa deployment where the attacker may not be
continuously present.  Nevertheless, it represents a formal gap that should be
closed before deployment.

The fix is a one-line firmware change: replace
\texttt{0x00000000} with \texttt{img\_seq} (a monotonically incrementing
\texttt{uint32} counter) in the \texttt{Sender\_Init} HMAC call.  This
modification brings the firmware into conformance with the Tamarin model and
closes the F3 gap.

\subsection{Future Work}
\label{sec:future}

\textbf{E20 --- FC-HMAC verifiability under packet loss (3$\times$ EoRa PI).}
A three-node experiment (sender, relay/dropper, receiver) will evaluate chain
verifiability at 10\% and 20\% packet drop rates, measuring the verifiable
prefix length and comparing FC-HMAC (chained) vs.\ per-packet HMAC (Ryczek
model).  The E20 firmware adds the \texttt{img\_seq} fix above, closing the F3
gap on hardware.

\textbf{E21 --- CPU/energy tradeoff measurement.}
A controlled comparison of \texttt{ENC\_NONE} / \texttt{ENC\_DC\_ONLY} /
\texttt{ENC\_FULL} on the same 2{,}046-byte JPEG, measuring both CPU
microseconds and, if a Joulescope power meter is available, milliwatt-hours per
image per mode.  The \texttt{ENC\_FULL} CPU estimate of $\approx$11{,}600~$\mu$s
used in this paper is derived from AES-128 software benchmarks on ESP32-S3;
E21 will provide a direct hardware measurement.

\textbf{Phase~3 live camera validation.}  The EoRa PI GPIO bank supports DVP
camera interface; attaching an OV2640 module will enable end-to-end experiments
with real captured images, validating that the stride-based scheme disrupts
natural (non-synthetic) JPEG images equally effectively.  Natural JPEG images
have more varied Huffman table layouts than synthetic test images; it is
important to confirm that the stride-at-byte-0 property holds across diverse
JPEG generators (e.g., libjpeg vs.\ ESP32-CAM's JPEG encoder).

\textbf{ROI-only partial transmission for true airtime reduction.}
As established in Section~\ref{sec:efficiency}, selective encryption alone does
not reduce airtime.  A future variant that identifies DCT-sensitive macroblocks
and \emph{omits} non-ROI fragments from transmission could reduce fragment count
to $\approx$24\% of the full image (2 of 9 fragments for a 200$\times$150 image
with a small ROI), cutting airtime from 14.9~s to $\approx$3.6~s.  This
requires the JPEG parser infrastructure that Phase~1 deliberately avoids, and
introduces partial-image reconstruction at the receiver.

\textbf{Per-session key derivation.}  The current pre-shared key model requires
manual key provisioning, which limits scalability.  A practical deployment
should use HKDF-derived per-session keys: $K_s = \mathrm{HKDF}(K_{\text{master}},
\texttt{img\_seq} \;\|\; \texttt{node\_id})$.  This eliminates the
\texttt{img\_seq} replay risk entirely (F3 gap) without requiring per-session
out-of-band key exchange.

\textbf{Multi-node LoRa mesh evaluation.}  The E19 and E20 experiments use
point-to-point LoRa links.  A LoRa mesh deployment with 5--10 nodes introduces
routing-layer fragmentation (LoRaMesher distance-vector routing) in addition to
application-layer JPEG fragmentation.  FC-HMAC must be evaluated in this
multi-hop context, where a relay node processes and re-transmits fragments and
may be compromised.  The formal Tamarin model currently assumes direct sender-to-receiver
communication; extending it to a relay model is planned for future work.

%% -----------------------------------------------------------------------
\section{Conclusion}
\label{sec:conclusion}

We presented Fragment-Chained HMAC (FC-HMAC), a lightweight authenticated
selective encryption scheme for JPEG image transmission over LoRa mesh networks.
FC-HMAC chains 8-byte HMAC-SHA256 tags across fragments and combines them with
a stride-based selective encryption that encrypts 24.2\% of image bytes.

The scheme achieves four properties simultaneously:
\begin{enumerate}
  \item \textbf{Visual obfuscation:} JPEG decode fails under both naive and
    format-aware attacker models (E19-PSNR evaluation, format-level disruption
    confirmed), while the CPU cost is 76\% lower than full encryption
    (2{,}807~$\mu$s vs.\ $\approx$11{,}600~$\mu$s for 2{,}046 bytes on ESP32-S3).
  \item \textbf{Airtime neutrality:} all three encryption modes (ENC\_NONE,
    ENC\_DC\_ONLY, ENC\_FULL) transmit 9 identical-length fragments at 14.9~s
    total airtime; the efficiency gain is strictly in CPU, not duty cycle.
  \item \textbf{Formal security:} four FC-HMAC lemmas proved in 1.53~s by
    Tamarin Prover under a Dolev-Yao adversary model, including chain-gap
    detectability (F2) and no-replay (F3) --- properties that per-packet
    independent HMAC cannot provide.
  \item \textbf{Hardware validation:} 72/72 fragments pass chain verification
    (100\% chain-ok, 100\% PDR) on EoRa PI hardware in a 120-second experiment
    at 922.0~MHz SF9/BW125, with 3.3\% per-packet HMAC overhead.
\end{enumerate}

FC-HMAC addresses the ordering and cross-session replay vulnerabilities of
per-packet independent HMAC at the cost of one 8-byte tag per fragment --- a
cost well within the 255-byte LoRa payload budget.  Future work will close the
F3 firmware gap by adding image sequence numbers to the chain root, validate
verifiability under 10--20\% packet loss in a three-node relay experiment,
and extend hardware measurements to live camera captures.

%% -----------------------------------------------------------------------
\bibliographystyle{IEEEtran}
\bibliography{references}

\end{document}
